\documentclass[12pt]{article}
\usepackage[utf8]{inputenc}
\usepackage[T1]{fontenc}
\usepackage{hyperref}
\title{Tasmanian, Port Sorell Language Resources for K-12 Curriculum Development}
\author{Compiled by Claude Sonnet 5.5 (claude-sonnet-5-5) for Andr\'as Kornai}
\date{}
\begin{document}
\maketitle

\section{Introduction}
Port Sorell is an extinct Tasmanian language of the north coast of Tasmania. It is known only from word lists, including handwritten ones by Charles Robinson held in the Mitchell Library. The Tasmanian Aboriginal community has revived a composite language, palawa kani, built from reconstructed vocabulary of the original languages. My searches found nothing specific to Port Sorell for schools.

\section{Primary Educational Resources}
I found no Port Sorell textbooks or primers. The palawa kani language program is the relevant community effort, but the Tasmanian Aboriginal Centre (TAC) restricts dictionaries and other learning materials to Aboriginal people. Others must request permission. Place names are stated to be freely available for public use.

\section{Linguistic and Reference Materials}
\begin{itemize}
\item Wikipedia, Port Sorell language: \url{https://en.wikipedia.org/wiki/Port_Sorell_language}
\item Wikipedia, palawa kani: \url{https://en.wikipedia.org/wiki/Palawa_kani}
\item The Oxford Guide to Australian Languages, chapter: \url{https://academic.oup.com/book/46706/chapter/410243287}
\item Endangered Languages Project, language information by source: \url{https://endangeredlanguages.com/node/111966}
\end{itemize}

\section{Religious and Community Resources}
\begin{itemize}
\item Tasmanian Aboriginal Centre, ``The policy and protocol for use of palawa kani Aboriginal language'': \url{https://tacinc.com.au/the-policy-and-protocol-for-use-of-palawa-kani-aboriginal-language/}
\end{itemize}

\section{Digital and Audio Resources}
\begin{itemize}
\item Google Arts and Culture story, ``palawa Country: Awakening Culture, Language and Landscape'': \url{https://artsandculture.google.com/story/RQWhq6xnl5GOfQ}
\end{itemize}

\section{Curriculum Development Context}
Any use of Tasmanian language material in schools should be led by the TAC and the Aboriginal community, and should follow the TAC protocol above. Material should be treated as heritage and revitalisation content, and no reconstructed vocabulary should be presented as Port Sorell speech without the community's agreement. The next step is to ask the TAC about the palawa kani Language Program, and to check the AIATSIS database entry (code T13) for Port Sorell.

\end{document}